% ch09_discussion.tex — Dissertation Chapter 9: Discussion and Future Work
% Blair Dupre, PhD Dissertation, University of North Dakota

\chapter{Discussion and Future Work}
\label{ch:discussion}

\section{Cross-Paper Synthesis}
\label{sec:synthesis}

The six papers comprising this dissertation form a vertically integrated computational pipeline, where each layer depends on the one below and adds value that the layers above require. Paper~1 (Chapter~\ref{ch:paper1}) establishes \textit{where the patient is now} through NSD-ISS stage classification. Paper~2 (Chapter~\ref{ch:paper2}) ensures that the input data supporting that classification is complete and trustworthy, even when clinical measurements are missing. Paper~3 (Chapter~\ref{ch:paper3}) extends the question temporally: \textit{where is the patient going, and when?} Paper~4 (Chapter~\ref{ch:paper4}) wraps every prediction in calibrated uncertainty, ensuring that clinicians receive not just point estimates but honest assessments of predictive confidence. Paper~5 (Chapter~\ref{ch:paper5}) tests whether these predictions remain valid when deployed on genuinely future patients. Paper~6 (Chapter~\ref{ch:paper6}) integrates all components into a unified system that answers the three clinical questions in a single interface.

This vertical integration reflects a deliberate architectural decision. Remove Paper~2 (imputation) and Paper~3's survival predictions become noisier because missing features degrade model input quality. Remove Paper~4 (conformal bands) and Paper~6's clinical report has no error bars---predictions without uncertainty cannot be used responsibly in clinical practice. Remove Paper~5 (temporal validation) and the argument for research-setting decision-support readiness collapses. The pipeline is designed so that each component is necessary and no component is sufficient alone.

The three-question clinical framework---\textit{What stage? When transition? How confident?}---emerged organically from the structure of the NSD-ISS staging system itself. The biological anchors (SAA and DaT-SPECT) define a patient's current position in disease space; the longitudinal trajectory through that space defines prognosis; and the irreducible heterogeneity of PD progression demands that every prediction be accompanied by a calibrated measure of uncertainty. The key results across the pipeline are: CatBoost AUC 0.981 for binary classification (Paper~1), GIMIN RMSE 107.7 with 22\% improvement over MissForest (Paper~2), Dynamic-DeepHit and Graph-DT both achieving C-td 0.926 for transition timing (Paper~3), IPCW conformal bands with 91.1\% marginal coverage at 95\% confidence (Paper~4), temporal validation C-td above 0.85 through three expanding windows (Paper~5), and end-to-end pipeline latency under two seconds per patient (Paper~6).


\section{Trees Versus Deep Learning on Clinical Data}
\label{sec:trees-vs-dl}

A recurring empirical finding across this dissertation is the dominance of gradient-boosted trees over deep learning on cross-sectional tabular clinical data, consistent with the comprehensive benchmarks of Grinsztajn et al.~\cite{grinsztajn2022} across 45 datasets and the analysis of Shwartz-Ziv and Armon~\cite{shwartzziv2022}. In Paper~1, CatBoost achieved AUC 0.981 for binary NSD-ISS classification while the Enhanced Multimodal GAT reached only 0.825---a gap of 15.6 percentage points. In Paper~2, MissForest and MICE outperformed all three deep learning imputation baselines (GAIN, SAITS, MIWAE) at every mask fraction, with the deep learning methods achieving RMSE values 1.5 to 2.4 times higher than the classical tree-based alternatives.

The explanation lies in the nature of clinical tabular data. PPMI features span wildly heterogeneous scales (clinical scores from 0--50, genetic risk scores near 50,000, binary indicators), mix continuous, categorical, and ordinal types, and exhibit complex non-linear interactions that gradient-boosted trees handle natively through recursive partitioning. Neural networks, by contrast, require careful normalization, are sensitive to feature scaling, and lack the rotational invariance that makes trees robust to arbitrary feature transformations~\cite{grinsztajn2022}.

However, this finding does not generalize to all tasks in the dissertation. Paper~3 requires \textit{sequential} modeling of variable-length visit histories and \textit{relational} modeling of patient similarity graphs---capabilities that tree-based methods fundamentally lack. CatBoost cannot process a sequence of twelve longitudinal visits or propagate information across a patient similarity network with 27,780 edges. The GRU component of Dynamic-DeepHit handles temporal dynamics; the GAT component of the Graph-Informed Digital Twin handles graph structure. This is not contradictory but rather appropriate tool selection: trees excel on cross-sectional tabular prediction, while deep learning excels when the data possesses temporal or relational structure that inductive biases can exploit.

The GIMIN architecture in Paper~2 represents an intermediate case. As a graph neural network operating on tabular features, GIMIN achieves RMSE 107.7---beating MissForest (137.3) by 22\%. The graph structure provides GIMIN with inter-patient relational information that tree-based imputation methods cannot access: the insight that a patient's missing cognitive score should be informed not just by their own observed features but by the observed cognitive scores of clinically similar patients in the graph neighborhood. This relational advantage is sufficient to overcome the general tabular disadvantage of neural networks, producing the strongest imputation results across the entire 12-model benchmark.


\section{The Imputation-Prediction Alignment Problem}
\label{sec:paradox}

Paper~2 reveals what may be the most conceptually important finding of the dissertation: the imputation-utility paradox. Stage-conditioned GIMIN variants achieve slightly \textit{worse} aggregate RMSE than the vanilla GIMIN (113.8 versus 107.7 at mask fraction 0.1), yet they produce \textit{better} downstream classification accuracy on the clinically most important staging targets (+2.9\% balanced accuracy for binary, +3.1\% for three-class). This paradox resolves when one examines per-stage RMSE: stage conditioning reallocates imputation capacity from the majority class (Stage~0, comprising 64.4\% of patients) to minority stages (Stages~1, 2B, and 4). Minority-stage RMSE improves by 6.4 to 24.3 points across mask fractions; Stage~0 RMSE degrades by up to 7.9 points.

This finding has broad implications beyond PD imputation. In any clinical dataset with severe class imbalance, optimizing for aggregate imputation error will systematically favor majority-class accuracy at the expense of the rare cases that are often the most clinically important. Balanced accuracy, which weights each class equally, reveals the true downstream impact. The lesson is that imputation and prediction serve different masters: RMSE rewards accuracy on common cases, while clinical utility rewards accuracy on the cases where correct classification matters most. Stage conditioning aligns the imputation objective with the downstream prediction objective by explicitly encoding the task-relevant structure of the problem.

We note the discovery order honestly: the minority-stage-advantage interpretation of the imputation-utility paradox was derived \emph{post~hoc} via the per-stage RMSE analysis, after we had already observed that stage-conditioned GIMIN variants improved downstream classification accuracy despite slightly worse aggregate RMSE. We did not pre-register a minority-stage-reallocation hypothesis before running the Paper~2 benchmark. The per-stage analysis should therefore be read as a \emph{consistency check on an observed pattern}, not as confirmation of a prior prediction. The mechanistic interpretation---capacity reallocation to rare stages---remains the best explanation consistent with the per-stage RMSE breakdown, but independent replication on a cohort with a different NSD-ISS stage prevalence is required to establish that the effect is robust to the specific class imbalance of PPMI rather than an artifact of it.

This pattern recurs throughout the dissertation. In Paper~3, the Graph-Informed Digital Twin matches Dynamic-DeepHit on mean C-td (0.926 versus 0.926) and shows numerically lower fold-to-fold standard deviation (0.013 versus 0.018), but the difference is not statistically significant (paired $t$-test $p=0.345$, Wilcoxon $p=0.312$). We note this as a trend suggesting potentially improved stability without overclaiming a variance advantage. In Paper~4, the IPCW-weighted conformal method achieves lower marginal coverage than naive conformal (82\% versus 90\% at 90\% confidence) but produces bands 2.6 times narrower---because it sacrifices coverage on trivially zero-probability cause-time combinations to achieve informative bands on clinically relevant transitions. In each case, the intermediate metric (aggregate RMSE, mean C-td, marginal coverage) tells one story, while the end-task metric (balanced accuracy, variance, band width) tells a different and more clinically relevant one.


\section{Limitations}
\label{sec:limitations}

Several limitations temper the conclusions of this dissertation and define the boundaries of its applicability. Before enumerating them, we note one methodological safeguard that is easy for a reader to overlook: every cross-validation split used in Papers~1--6 is \emph{patient-blocked}. Specifically, \texttt{StratifiedKFold} is applied to the list of unique PATNOs (not to the episode or visit list), and all episodes belonging to a given patient are then routed to the same fold via set-membership filtering (see \texttt{src/giman\_pipeline/paper3/dynamic\_deephit.py} and \texttt{graph\_digital\_twin.py} fold-assignment routines). Paper~4's conformal calibration/evaluation split operates at the patient level for the same exchangeability reason. This eliminates cross-fold patient leakage as a source of optimistic bias in the reported C-td and conformal-coverage numbers.

\textbf{Class imbalance.} The NSD-ISS stage distribution in PPMI is severely imbalanced: Stage~0 comprises 64.4\% of patients, while Stage~4 comprises only 0.8\%. This imbalance propagates through every downstream analysis. Paper~1's full ordinal classification achieves the lowest balanced accuracy (0.660), driven by near-random performance on Stage~4. Paper~2's aggregate RMSE is dominated by Stage~0 imputation quality. Paper~3's per-transition C-td varies from 0.856 (transitions to Stage~0) to 0.944 (transitions to Stage~4), reflecting the uneven event counts across transition types. While techniques such as class-weighted loss functions, balanced accuracy metrics, and stage-conditioned architectures partially mitigate this imbalance, the fundamental limitation is small sample sizes for rare stages.

\textbf{SAA coverage.} Only 12.6\% of PPMI patients (277 of 2,201) have SAA results available, meaning the S~anchor---the molecular foundation of NSD-ISS staging---is missing for 87.4\% of the cohort. Staging for these patients relies entirely on the D~anchor (DaT-SPECT, 97.1\% coverage). This introduces a systematic bias: patients with only D~anchor data may be misclassified between Stage~0 and Stage~1 if they are S+ but D$-$. The imputation framework of Paper~2 partially addresses this by enabling imputation of missing SAA status, but the fundamental limitation of sparse S~anchor coverage remains.

\textbf{Single-cohort training (temporal vs.\ external generalization are distinct).} All survival models (Papers~3--6) are trained exclusively on PPMI data, and it is important to distinguish two very different kinds of generalization claim that this dissertation makes. Paper~5's temporal validation tests robustness to \emph{future PPMI patients}---an intra-cohort, within-distribution shift captured by the W1--W4 expanding windows, including the W4 stress test. This is generalization \emph{over time within a single cohort}. Paper~1's external validation, in contrast, tests robustness across \emph{different cohorts} (BioFIND, PDBP, HBS), which introduces ascertainment, protocol, and population differences that are categorically distinct from temporal drift. Paper~5's temporal validation should not be read as a substitute for external cohort validation, and Paper~1's external validation should not be read as a substitute for temporal stability. No external cohort has longitudinal NSD-ISS staging with sufficient follow-up for transition modeling, so Papers~3--6's survival results remain untested against true domain shift; this is a real limitation, not something temporal validation alone can close.

\textbf{PPMI cohort characteristics.} PPMI is a research cohort with standardized assessment protocols, relatively complete data collection, and a participant population that skews toward higher education and socioeconomic status compared to the general PD population. Clinical data from community neurology practices will be noisier, more incomplete, and drawn from a more diverse population. The degree to which models trained on PPMI will generalize to routine clinical settings remains an open empirical question.

\textbf{Medication confound.} Espay et al.~\cite{espay2025} have critiqued the NSD-ISS staging system for its reliance on functional assessments (UPDRS-III, ADL scales) that are sensitive to symptomatic medication effects. A patient on optimized levodopa therapy may appear functionally less impaired (Stage~2B) than their underlying biology warrants (perhaps Stage~3 or 4 without medication). The 39.1\% backward transition rate observed in Paper~3 is at least partially attributable to medication initiation or optimization rather than true biological regression. The NSD-ISS staging system conflates symptomatic benefit with biological stage, and this conflation propagates through every model in the pipeline.

\textbf{Domain shift in external validation.} Paper~1's primary external clinical result is the NSD-positive sub-staging AUC of 0.900 on BioFIND, which measures the clinically meaningful task of discriminating sub-stages \emph{among diagnosed PD patients}. The binary classification AUC on BioFIND (0.637) is \emph{not} a clinical performance number and should not be interpreted as one: it is best read as a data-quality diagnostic that surfaces PPMI's inclusion of healthy controls in Stage~0. Because PPMI's Stage~0 blends healthy controls with biologically early patients while BioFIND is a PD-only cohort, a model trained on PPMI binary labels effectively learns an HC-vs-PD decision boundary rather than an S$+$-vs-S$-$ boundary, and that learned boundary collapses under the BioFIND distribution. Three-class AUC (0.703 external) sits between the two, consistent with partial but not total confounding. We therefore report NSD-positive sub-staging as the external clinical result, the three-class AUC as an intermediate robustness check, and the binary external AUC as a diagnostic of PPMI's cohort composition rather than a statement about downstream clinical performance.


\section{Future Work: Toward Mechanistic Digital Twins}
\label{sec:future}

The GIMAN framework established in this dissertation is fundamentally \textit{correlational}: it learns statistical associations between observed features and future stage transitions. It cannot answer the interventional question that matters most for treatment: ``What if this patient starts prasinezumab now versus in twelve months?'' A mechanistic digital twin would change this entirely by encoding the causal biological chain---alpha-synuclein aggregation driving neuronal death, which reduces striatal dopamine, which manifests as declining DaT-SPECT signal and worsening motor function. By encoding mechanism, the model becomes interventional, capable of simulating outcomes under treatments that alter specific biological processes.

The transition from correlational to mechanistic modeling requires five coupled ordinary differential equation (ODE) modules, each grounded in established disease biology and calibrated from data that already exist within AMP-PD. A comprehensive technical treatment of the mathematical formulations, data sources, calibration strategies, and validation criteria for each module is provided in Appendix~\ref{app:mechtwin}; this section summarizes the key requirements and identifies which components can be built now versus those that require future data.

The \textbf{alpha-synuclein aggregation module} encodes nucleation-elongation kinetics~\cite{knowles2009,dear2020,meisl2014,cohen2012}: monomeric alpha-synuclein misfolds into oligomeric nuclei that template fibril formation, governed by production, nucleation, elongation, fragmentation, and clearance rate constants. In-vitro kinetic parameters from ThT fluorescence assays~\cite{knowles2009,meisl2014} provide informative priors, while approximately 1,200 PPMI patients with CSF total alpha-synuclein measurements (1--2 timepoints) and 277 with binary SAA results provide sparse but usable calibration data. Anti-alpha-synuclein antibodies such as prasinezumab~\cite{pagano2022,pagano2024} enter these equations by reducing the elongation rate constant, enabling simulation of treatment effects without requiring treatment data in the training set. Next-generation quantitative SAA assays, which provide continuous fluorescence readout proportional to fibril mass~\cite{shahnawaz2023}, are expected within two to three years and would transform the calibration landscape for this module.

The \textbf{dopaminergic neuron death module} models the progressive loss of substantia nigra neurons driven by oligomeric toxicity, neuroinflammation, and age-related attrition, with DaT-SPECT SBR serving as a noisy proxy for surviving neuron count. This is the \textit{most immediately buildable} module: PPMI provides serial DaT-SPECT for approximately 2,000 patients with 2--5 scans over 5--7 years~\cite{marek2011}, enabling per-patient Bayesian estimation of the death rate constant $k_\text{death}$. Longitudinal neuromelanin MRI~\cite{xing2021} for a subset of 300--500 patients provides a complementary, more direct proxy for surviving neuron count. Genetic status (LRRK2, GBA)~\cite{simuni2020lrrk2} modulates decline rates and can be incorporated as patient-specific rate modifiers.

The \textbf{Lewy body propagation module} implements network diffusion across the structural connectome following the Braak hypothesis~\cite{braak2003}, extending the mathematical framework of Raj et al.~\cite{raj2012,raj2019} and Fornari et al.~\cite{fornari2019} for prion-like cell-to-cell transmission of misfolded alpha-synuclein~\cite{brundin2010}. The alpha-synuclein origin and connectome (SOC) model~\cite{horsager2022} proposes that the origin site of pathology determines clinical phenotype, motivating patient-specific seeding location. DTI tractography from approximately 400 PPMI patients provides individual structural connectomes; the Human Connectome Project (HCP) population-average connectome serves as a reference for the remaining patients. FreeSurfer regional volumes ($\sim$1,700 PPMI patients, 1--3 serial MRI scans) provide downstream observables of regional pathology density.

The \textbf{pharmacokinetic/pharmacodynamic module} captures levodopa absorption, distribution, metabolism, and the Hill-type dose-response relationship between synaptic dopamine and motor function~\cite{ursino2020}. This module benefits from the richest data: PPMI medication records ($\sim$59,000 concomitant medication entries) linked to serial UPDRS-III assessments across 16,699 visits provide input-output pairs for calibrating the dose-response relationship. Population PK parameters ($k_a \approx 1.5$~hr$^{-1}$, $k_\text{el} \approx 0.7$~hr$^{-1}$) are available from published models; patient-specific PD parameters (EC$_{50}$, Hill coefficient) are calibratable from levodopa equivalent daily dose (LEDD) versus UPDRS-III response curves. Gut microbiome modulation of levodopa bioavailability~\cite{vankessel2019} introduces a source of inter-individual variability that future iterations should capture.

Finally, the \textbf{functional impairment module} maps the biological state vector to clinically observable NSD-ISS stage---this module is essentially already built as Paper~1's CatBoost classifier, with the 12-feature clinical-only model (AUC 0.900 for NSD-positive sub-staging) providing the most natural interface since its inputs map directly to mechanistic module outputs.

GIMAN provides the foundation upon which the mechanistic framework builds, and this integration distinguishes the proposed approach from standalone mechanistic models~\cite{bjornsson2020,laubenbacher2024}. Paper~1's staging classifier serves as the functional impairment module and as a validation oracle for mechanistic model outputs. Paper~2's imputation framework handles missing serial biomarker data needed for parameter calibration---for instance, imputing missing CSF alpha-synuclein values for patients who lack lumbar puncture data. Paper~3's empirical transition rates and Markov sojourn times (Stage~0: 13.3~years, Stage~2B: 0.68~years, Stage~3: 1.85~years, Stage~4: 1.42~years) serve as direct calibration targets for the neuron death rate constant $k_\text{death}$: the mechanistic model's predicted $N(t)$ trajectory must produce transition times consistent with these empirical rates. Paper~4's conformal bands provide a rigorous validation criterion: mechanistic predictions that consistently fall outside GIMAN's IPCW conformal bands (width 0.037 at 95\% confidence) signal parameter miscalibration. The patient similarity graph from Paper~3 (1,900 nodes, 27,780 edges) provides graph-regularized priors for Bayesian parameter estimation~\cite{villaverde2022}, where a patient's rate constants are informed by the already-calibrated parameters of their graph neighbors---an approach that partially addresses the fundamental identifiability challenge of calibrating a multi-parameter ODE system from sparse clinical observations.

\begin{figure}[t]
\centering
\begin{tikzpicture}[
    node distance=1.0cm,
    databox/.style={rectangle, draw, fill=green!8, text width=4.5cm, minimum height=1.0cm, align=center, font=\small, rounded corners=2pt},
    mechbox/.style={rectangle, draw, fill=orange!10, text width=4.5cm, minimum height=1.0cm, align=center, font=\small, rounded corners=2pt},
    readybox/.style={rectangle, draw, fill=orange!25, text width=4.5cm, minimum height=1.0cm, align=center, font=\small, rounded corners=2pt, line width=0.8pt},
    bridgebox/.style={rectangle, draw, fill=purple!8, text width=10.2cm, minimum height=0.8cm, align=center, font=\small, rounded corners=2pt},
    arrow/.style={-stealth, thick, draw=gray!70},
    coupl/.style={-stealth, thick, draw=red!50, dashed},
    brace/.style={decorate, decoration={brace, amplitude=6pt, raise=2pt}, thick}
]
% Left column: GIMAN (data-driven)
\node[databox] (p1) {Paper~1: Stage Classification\\CatBoost AUC 0.981};
\node[databox, below=0.4cm of p1] (p2) {Paper~2: GIMIN Imputation\\RMSE 107.7};
\node[databox, below=0.4cm of p2] (p3) {Paper~3: Transition Timing\\C-td 0.926};
\node[databox, below=0.4cm of p3] (p4) {Paper~4: Conformal Bands\\91.1\% coverage};
\node[databox, below=0.4cm of p4] (p5) {Paper~5: Temporal Validation\\C-td $>$ 0.85};

% Right column: Mechanistic (future) — darker = more data-ready
\node[databox, right=1.2cm of p1] (m5) {Module 2e: Functional Mapping\\{\scriptsize Already built (Paper~1)}};
\node[mechbox, right=1.2cm of p2] (m1) {Module 2a: $\alpha$-Syn ODEs\\{\scriptsize CSF: $\sim$1,200 pts, SAA: 277}};
\node[readybox, right=1.2cm of p3] (m2) {Module 2b: Neuron Death\\{\scriptsize DaT-SPECT: $\sim$2,000 pts}};
\node[mechbox, right=1.2cm of p4] (m3) {Module 2c: Connectome Diffusion\\{\scriptsize DTI: $\sim$400; FreeSurfer: $\sim$1,700}};
\node[readybox, right=1.2cm of p5] (m4) {Module 2d: Levodopa PK/PD\\{\scriptsize Meds: 59K records; UPDRS: 16.7K}};

% Horizontal arrows (GIMAN provides foundation)
\draw[arrow, dashed] (p1.east) -- (m5.west) node[midway, above, font=\scriptsize] {serves as};
\draw[arrow, dashed] (p2.east) -- (m1.west) node[midway, above, font=\scriptsize] {imputes for};
\draw[arrow, dashed] (p3.east) -- (m2.west) node[midway, above, font=\scriptsize] {calibrates};
\draw[arrow, dashed] (p4.east) -- (m3.west) node[midway, above, font=\scriptsize] {validates};
\draw[arrow, dashed] (p5.east) -- (m4.west) node[midway, above, font=\scriptsize] {temporally validates};

% Inter-module coupling arrows (right column)
\draw[coupl] (m1.south) -- (m2.north) node[midway, right, font=\scriptsize, text=red!60] {$O(t)$};
\draw[coupl] (m2.south) -- (m3.north) node[midway, right, font=\scriptsize, text=red!60] {$N(t)$};
\draw[coupl] (m1.south east) to[out=-30,in=30] (m3.north east) node[midway, right, font=\scriptsize, text=red!60] {};

% Braces and labels
\draw[brace] ([xshift=-0.3cm]p1.north west) -- ([xshift=-0.3cm]p5.south west) node[midway, left=0.4cm, align=center, font=\small\sffamily] {Data-Driven\\(This Dissertation)};
\draw[brace] ([xshift=0.3cm]m4.south east) -- ([xshift=0.3cm]m5.north east) node[midway, right=0.4cm, align=center, font=\small\sffamily] {Mechanistic\\(Future Work)};

% Bridge box at bottom
\node[bridgebox, below=0.8cm of $(p5.south)!0.5!(m4.south)$] (bridge) {MindMend Biosensor: Continuous $\alpha$-synuclein $\rightarrow$ closed-loop state estimation};
\draw[arrow] (p5.south) -- ++(0,-0.35) -| (bridge.north west);
\draw[arrow] (m4.south) -- ++(0,-0.35) -| (bridge.north east);

\end{tikzpicture}
\caption{Transition from data-driven (GIMAN, this dissertation) to mechanistic digital twin (future work). Each GIMAN paper provides a specific component that the corresponding mechanistic module builds upon. Darker orange boxes indicate modules with sufficient PPMI data for immediate calibration; lighter boxes require future data. Red dashed arrows show inter-module coupling: oligomer concentration $O(t)$ drives neuron death, surviving neuron count $N(t)$ modulates drug response, and fibrillar $\alpha$-synuclein feeds connectome propagation. The MindMend wearable biosensor bridges both paradigms.}
\label{fig:mech-transition}
\end{figure}

Table~\ref{tab:data-gaps-ch9} summarizes the critical data gaps constraining the timeline for mechanistic model development, their severity, and available mitigations. The pragmatic development path follows a data-availability-first strategy: begin with the neuron death module (Module~2b, richest longitudinal data in serial DaT-SPECT) and the PK/PD module (Module~2d, extensive medication + UPDRS records), then progressively add upstream biological modules as quantitative biomarker assays become available. The functional impairment module (Module~2e) is already built as Paper~1's CatBoost classifier.

\begin{table}[t]
\centering
\caption{Critical data gaps for mechanistic model development.}
\label{tab:data-gaps-ch9}
\small
\begin{tabular}{p{3.2cm}lp{3.5cm}l}
\toprule
\textbf{Data Need} & \textbf{Severity} & \textbf{Mitigation} & \textbf{Timeline} \\
\midrule
Quantitative CSF $\alpha$-syn species & High & Literature priors + total CSF $\alpha$-syn & 2--3~yr (qSAA) \\
Serial CSF (>2 pts/patient) & High & Graph-regularized Bayesian priors & Ongoing \\
Individual DTI connectome & Medium & HCP population average & Available now \\
Plasma levodopa PK & Medium & Population PK from literature & Ancillary study \\
Continuous $\alpha$-syn monitoring & High & MindMend biosensor & 3--5~yr (TRL~6) \\
\bottomrule
\end{tabular}
\end{table}

The MindMend wearable biosensor (Patent US20250334570A1) addresses the most critical data gap---continuous alpha-synuclein monitoring---through interstitial fluid measurement using graphene oxide functionalized electrochemical sensors with anti-alpha-synuclein aptamers. At a temporal resolution of one to five minute intervals yielding 300--1,440 data points per day, the biosensor would transform the aggregation module from open-loop prediction (calibrated from one to two CSF samples per patient lifetime) to closed-loop state estimation via particle filter or ensemble Kalman filter updates. The technology is currently at TRL~2 (concept formulated, patent published); a realistic timeline to first-in-human pilot is 30--42 months, with closed-loop digital twin demonstration at 42--60 months.

The recommended computational infrastructure centers on Julia's \texttt{DifferentialEquations.jl} ecosystem~\cite{rackauckas2017,rackauckas2020}, which provides stiff-aware ODE solvers (CVODE\_BDF for the multi-timescale coupled system), automatic differentiation for sensitivity analysis, and GPU acceleration for population-level calibration. Bayesian parameter estimation uses Turing.jl (NUTS/HMC), with structural identifiability verified via \texttt{StructuralIdentifiability.jl} before calibration---a prerequisite given the severe under-determination of the aggregation module (9 parameters, 1--2 observations). The physics-informed machine learning paradigm~\cite{karniadakis2021} and universal differential equations framework~\cite{rackauckas2020} provide the mathematical foundation for hybrid mechanistic-ML models that combine ODE constraints with neural network residuals.

A phased five-year development plan proceeds as follows. Phase~1 (months 1--6) implements the coupled alpha-synuclein aggregation and neuron death ODE system in Julia with synthetic data validation and identifiability analysis. Phase~2 (months 4--12) performs Bayesian calibration of patient-specific neuron death rates from serial DaT-SPECT using graph-informed priors from GIMAN's patient similarity network, producing posterior distributions for approximately 2,000 PPMI patients---this phase is immediately actionable because DaT-SPECT data with 3--4 serial timepoints per patient already exist. Phase~3 (months 10--24) implements the connectome diffusion propagation model~\cite{raj2012,raj2019,fornari2019} calibrated from FreeSurfer volumetric data, with the HCP population-average connectome providing structural connectivity for patients without individual DTI. Phase~4 (months 12--24) integrates the levodopa PK/PD module calibrated from PPMI medication records and UPDRS trajectories. Phase~5 (months 20--30) performs head-to-head comparison of mechanistic predictions against GIMAN's data-driven predictions, with hybrid model development if mechanistic components add value---the Universal Differential Equations framework~\cite{rackauckas2020} enables embedding neural network residual terms within the ODE system to capture unmodeled biology. Phase~6 (months 24--60) integrates MindMend biosensor data for real-time closed-loop digital twin operation via ensemble Kalman filter~\cite{cheng2023}. The estimated total resource requirement is approximately \$770,000 over four years, well within the scope of an NIH R01 or R21 mechanism.

The envisioned clinical impact of the mechanistic framework is qualitatively different from GIMAN's correlational predictions. Where GIMAN reports ``this patient has a 42\% probability of transitioning to Stage~3 within 24 months,'' a mechanistic twin would report ``this patient's calibrated aggregation rate predicts Stage~3 arrival at month~16. If prasinezumab is initiated now, reducing the elongation rate by an estimated 25\%, predicted arrival shifts to month~28---a delay of 12 months.'' The mechanistic model provides counterfactual predictions with quantified treatment benefits, enabling truly personalized therapeutic decision-making.
